\chapter{Artist Objects}\label{chapter:artists}

Artist objects are the water we've been swimming in this whole time. Everything rendered on a plot is some kind of Artist object. Even the figure and axes objects belong to this base class. Why notice Artist objects now? Because we're leaving the well-worn path of line plots, histograms, and so forth, and some appreciation of our surroundings will help. It will also provide some comfort as you continue to learn more and encounter more technical documentation.

Chapter \ref{chapter:elements} introduced the distinction between primitive Artists, such as lines, text, and patches, and the containers that hold them. We do not need a deeper theory of Artist objects here. It is enough to recognize them when we create them and know how to add them to an axes object. The plots in the following chapter might look unusual, but they are mostly arrangements of lines, text, patches, and images.

\section{Adding Artist Objects}

Many axes methods quietly create Artist objects for us. For example, \code{ax.plot()} creates one or more \code{Line2D} objects, \code{ax.text()} creates a \code{Text} object, and \code{ax.scatter()} creates a collection. A call to \code{ax.bar()} creates a collection of rectangular patches. You usually do not need to know any of this when the usual plotting method gives you what you want.

But sometimes it is easier to create the Artist directly. You can create a circle with \code{plt.Circle()} and add it to an axes object with \code{ax.add\_patch()}. There are similar methods such as \code{add\_line()} and \code{add\_collection()}, along with the more general \code{add\_artist()}. I would use the more specific method when one exists. Besides making the type of object clearer, some of these methods also allow the axes object to account for the Artist when determining the data limits. Matplotlib's \link{https://matplotlib.org/stable/tutorials/artists.html}{Artist tutorial} goes into more detail.

Creating and adding an Artist does not finish the job. We still need to choose its coordinate system, its \code{zorder}, and whether it should be clipped by the axes. We might also set its color, transparency, line width, or size. These are not new considerations. We have already made the same choices through plotting methods. Creating the Artist directly just leaves more of those decisions in our hands.

\section{Building a Plot One Artist at a Time}

When I build a custom plot, I usually begin by creating the figure and axes objects and setting the axes limits or coordinate system. Then I create the required lines, patches, images, and text and add them to the axes object. The \code{zorder} can be used to arrange the layers. Only after seeing everything together will you know whether the spacing, clipping, and sizing work. This is one reason these plots often require some tinkering.

This is more laborious than calling \code{plot()}, but nothing fundamentally different is happening. We are still using Matplotlib's object-oriented interface, only one object at a time. The activity calendar, heatmaps, directed graph, and speedometer that follow look very different, but each is built by choosing a few Artist objects and placing them on an axes object.
